\documentclass[10pt,a4paper]{amsart}
\usepackage[margin=19mm]{geometry}
\usepackage{amsmath,amssymb,mathtools}
\usepackage[scr=rsfs,cal=euler]{mathalfa}
\usepackage{xfrac}
\usepackage[unicode,hidelinks,bookmarks=false]{hyperref}
\newcommand{\ie}{i.e.\ }
\newcommand{\cf}{cf.\ }
\newcommand{\resp}{respectively\ }
\newcommand{\Subsection}{\subsection}
\providecommand{\nobreakdash}{\nobreak\hbox{-}}
\setlength{\emergencystretch}{2em}
\multlinegap0pt
\usepackage{bm}
\usepackage{bbm}
\usepackage{dsfont}
\usepackage{xspace}
\usepackage{cancel}
\usepackage{mathtools}
\usepackage{amsthm}
\makeatletter
\@ifundefined{Proof}{\newtheorem{Proof}{Proof}}{}
\@ifundefined{Theorem}{\newtheorem{Theorem}{Theorem}}{}
\makeatother
\title[Indefinite Integration Operators Identities and their Approx]{Indefinite Integration Operators Identities and their Approximations}
\author{Frank Stenger}
\date{Source: arXiv:1809.05607v2 (CC BY 4.0)}
\begin{document}
\maketitle

\noindent\emph{Source and licence.} Indefinite Integration Operators Identities and their Approximations. Version arXiv:1809.05607v2 (CC BY 4.0), distributed under the Creative Commons Attribution 4.0 International licence, as recorded in the arXiv licence metadata for this exact version and shown on the abstract page https://arxiv.org/abs/1809.05607v2. Full rights notice: licenses/arxiv-1809.05607-CC-BY-4.0.txt.\par\smallskip

\noindent\textit{Editorial scope.} Counted unit: the Theorem whose source label is \texttt{theorem:2.5} in the article (source0.tex lines 421--433, source label \texttt{theorem:2.5}), delivered together with the article's own complete original proof of it (source0.tex lines 435--514, one \texttt{proof} environment of 80 lines). No mathematics is written, repaired or re-derived: statement and proof are the article's own text line for line. Required context retained uncounted: none. Every definition, display and label of those lines is retained unchanged. Omitted from this package and not redistributed: 1--420, 434--434, 515--1567. Nothing inside a retained excerpt is omitted or altered: every retained body is delivered exactly as the article's lines are written, so no omission receipt of any kind accompanies this package. Citations: the retained excerpts carry no citation command at all, so no bibliography entry of the article is transcribed and no reference is redistributed. Reference adaptation: none was needed and none was made; every reference inside the delivered unit resolves inside it, so no unresolved reference or citation is produced. Printed slips kept literal: no printed word, symbol, hypothesis or number of the counted statement or of its proof was altered. Layout and class: the article's own class is \texttt{article} with packages including amsmath, amssymb, amsthm, graphicx; the wrapper is the lane's pinned \texttt{amsart} editorial wrapper at 10pt on A4, which restates the theorem-like environments the retained text needs and adds the article's own macro definitions that the retained text needs (none), with extra wrapper packages (bm, bbm, dsfont, xspace, cancel, mathtools). The delivered numbering is the unit's own. Assets: no retained span contains a figure environment, a graphics-inclusion command, a PDF-page inclusion, a table or a TikZ picture; no image file is copied into this package, the package declares no asset dependency and no image file is redistributed with it. Licence: version arXiv:1809.05607v2 of this article is distributed under the Creative Commons Attribution 4.0 International licence, as recorded in the arXiv licence metadata for this exact version and shown on the abstract page https://arxiv.org/abs/1809.05607v2. A full rights notice is delivered at licenses/arxiv-1809.05607-CC-BY-4.0.txt. Characters: every retained excerpt is pure ASCII, so the delivered engine, XeLaTeX with its default Latin Modern text font, reports no missing character.
  \begin{Theorem}%2.5
    \label{theorem:2.5}
Let ${\mathcal J}^\mp$ be defined as above, and have support on $(0,b)
\subseteq {\mathbf R}^+$\,, let the functions $f^\mp$ of equation
(2.6) belong to ${\mathbf H}$\,.  Let $q^\mp$ be defined as in
(2.6).  Then
 
 \begin{equation}%2.10
 q^\mp = \widehat{f^\mp}\left(\mp i\,/{\mathcal J}^\mp\right)\,g\,.
\label{equation:2.10}
 \end{equation}

 \end{Theorem}

 \begin{Proof} The proof resembles that given in \S 4.6 of [Sp] and
   \S1.5.9 of [Sc] for the case of Laplace transforms. We consider
   only the case of $f^+$\,, since the proof for the case of $f^-$ is
   similar.     The proof makes use of the following:
 \begin{itemize}
   \item $f^+$ has compact support $[0,b]$\,, -- inasmuch
       extension to $(0,\beta)$\,, or to an infinite interval can be
       carried out via the usual well--known procedure of analysis,
       and moreover, our assumptions are consistent with this possibility;
   \item By Theorem 2.4, the spectrum of $i/{\mathcal J}^+$ lies in
     the closed second quadrant of the complex plane, ${\mathbf C}$\,.
   \item As can be shown via an easy to prove inequality -- see
       [Sp], \S4.6 -- namely, that $\left\|{\mathcal
         J}^\mp\right\|_p \leq \alpha$ for all $p \in [1,\infty]$\,,
       where $(0\,,\alpha) \subset {\mathbf R}^+$ is the support of
       ${\mathcal J}^\mp$\,;
   \item By inspection of (2.6), $\widehat{f^+}$ is
     analytic in the upper half of the complex plane, and
      $\widehat{f^-}$ is analytic in the lower half; and
   \item The well--known formulas (2.6), as well as  

   \begin{equation} %2.11
     \begin{array}{rcl}
          \displaystyle \left(\left({\mathcal J}^+\right)^n\,g\right)(x) & = & 
          \displaystyle \int_0^x 
          \frac{(x-t)^{n-1}}{(n-1)\,!} \, g(t) \, dt\,, \ \ \ n =
          1\,,2\,,\ldots\,,\\ 
           \\
          \widehat{f^\mp}(z)  & = & \displaystyle \mp
          \frac{1}{2\,\pi\,i} \int_{\mathbf R} 
          \frac{\widehat{f^\mp}(t)}{t - z} dt\,, \ \ \ \mp\,\Im z > 0\,.
     \end{array}
    \label{equation:2.11}
   \end{equation}
   
 \end{itemize}       

\noindent Hence, applying the above points, for $x \in (0,b)$\,, and with $I$
denoting the identity operator, we get

%2.12
\begin{equation}
  \begin{array}{rcl}
q^+(x) & = & \displaystyle \int_0^x f^+(x-\xi) \, g(\xi) \, d\xi \\
    && \\
    & = & \displaystyle \int_0^x \,\left(\frac{1}{2\,\pi} \, \int_{\mathbf R}
\widehat{f^+}(t) \, e^{-\,i\,t\,(x - \xi)}\,dt \,
\right)\,g(\xi)\,d\xi \\  
&& \\
& = & \displaystyle \frac{1}{2\,\pi} \,\int_{\mathbf R} \left(\int_0^x
\,\sum_{n=0}^\infty  
         \frac{(-i\,t\,(x-\xi))^n}{n\,!} \, 
         g(\xi)\,d\xi\right) \, \widehat{f^+}(t) \, dt \\ 
         && \\
& = & \displaystyle \left(\frac{1}{2\,\pi} \, \int_{\mathbf R}
         {\mathcal J}^+ \, 
         \sum_{n=0}^\infty \left(-\,i\,t\,{\mathcal J}^+\right)^n \,
         \widehat{f^+}(t) \, dt \, g\right)(x) \\ 
         && \\
& = & \displaystyle \left(\left(\frac{1}{2\,\pi\,i} \int_{\mathbf R} {\mathcal
           J}^+/\left(1 + i\,t\,{\mathcal J}^+\right) \,
         \widehat{f^+}(t) 
         \, dt\right) \, g \, \right)(x) \\
&& \\
     & = & \displaystyle \lim_{\delta \to 0^+}\left(\frac{1}{2\,\pi\,i}\,\int_{\mathbf R}
      \widehat{f^+}(t)/\left(t\,I - i/\left({\mathcal
          J}^++\delta\,I\right)\right) \, dt \, 
      g\right)(x) \\ 
      && \\
      & = & \displaystyle \lim_{\delta \to 0^+}
      \left(\widehat{f^+}\left(i/\left({\mathcal J}^+ +
      \delta\,I\right)\right)\,g\right)(x) = 
      \left(\widehat{f^+}\left(i/{\mathcal J}^+\right)\,g\right)(x)\,,
  \end{array}
      \label{equation:2.12}
\end{equation} 

\noindent Similarly, $q^-(x) = \left(\widehat{f^-}\left(-i/{\mathcal
  J}^-\right)\,g\right)(x)$\,.  
  \end{Proof} $\blacksquare$

\end{document}
